\section{Mechanism Design}

\label{sec:mechanism}
% ══════════════════════════════════════════════════════════════════════════════

\subsection{Stakeholders}

A BIB involves five stakeholder roles:

\begin{enumerate}
  \item \textbf{Investors} ($I$): Provide upfront capital by purchasing
    BIB tokens, seeking financial returns contingent on outcomes.
  \item \textbf{Implementers} ($M$): Conservation organizations that
    execute on-the-ground conservation activities.
  \item \textbf{Outcome payers} ($O$): Entities (governments,
    corporates, philanthropies) that commit to paying for verified
    outcomes.
  \item \textbf{Verifiers} ($V$): Decentralized oracle network that
    assesses outcome achievement through monitoring data.
  \item \textbf{Governance} ($G$): Zoo DAO, which sets standards,
    arbitrates disputes, and manages protocol parameters.
\end{enumerate}

\subsection{Outcome Metrics}

BIB outcomes are specified using the \textbf{SMART-B} framework
(Specific, Measurable, Attributable, Realistic, Time-bound,
Biodiversity-relevant):

\begin{table}[H]
\centering
\caption{Example outcome metrics for BIB project archetypes.}
\label{tab:metrics}
\begin{tabular}{p{3cm}p{4cm}p{4cm}}
\toprule
\textbf{Project type} & \textbf{Primary metric} & \textbf{Verification} \\
\midrule
Species recovery & Population count $\geq$ threshold & Camera traps + eDNA \\
Habitat protection & Forest area maintained (ha) & Satellite NDVI \\
Corridor creation & Connectivity index $\geq$ target & Movement tracking \\
Reef restoration & Coral cover (\%) increase & Underwater surveys \\
Invasive removal & Target species abundance decrease & Transect surveys \\
\bottomrule
\end{tabular}
\end{table}

\subsection{Payment Schedule}

BIBs use a milestone-based payment structure with $T$ verification
periods over the bond lifetime $[0, \tau]$:

\begin{equation}
  \text{Payment}_t = \begin{cases}
    \pi_t^{\text{base}} + \pi_t^{\text{bonus}} \cdot \mathbb{1}[O_t \geq O_t^*]
    & \text{if } O_t \geq O_t^{\min} \\
    0 & \text{otherwise}
  \end{cases}
\end{equation}

where $O_t$ is the measured outcome at period $t$, $O_t^{\min}$ is the
minimum threshold, $O_t^*$ is the target for full payment,
$\pi_t^{\text{base}}$ is the base payment, and $\pi_t^{\text{bonus}}$
is the performance bonus.

\subsection{Token Economics}

\paragraph{Issuance.} A BIB issuance creates $N$ fungible tokens,
each representing a fractional claim on the bond's payment stream:
\begin{equation}
  P_0 = \frac{C_{\text{total}} + F_{\text{admin}}}{N}
\end{equation}

\paragraph{Returns.} Token holders receive pro-rata shares of outcome
payments:
\begin{equation}
  \E[\text{Return}] = \frac{1}{P_0} \sum_{t=1}^{T}
  \frac{\pi_t^{\text{investor}}}{(1+r)^t} - 1
\end{equation}

\paragraph{Secondary market.} Tokens are transferable on decentralized
exchanges, with pricing reflecting updated outcome probability
assessments from oracle data.

\subsection{Capital Flow Architecture}

\begin{figure}[H]
\centering
\begin{tikzpicture}[
  node distance=2cm,
  box/.style={draw, rounded corners, minimum width=2.5cm,
    minimum height=1cm, align=center, font=\small},
  arrow/.style={-{Stealth[length=3mm]}, thick}
]
  \node[box] (investors) {Investors};
  \node[box, right=3cm of investors] (escrow) {Smart Contract\\Escrow};
  \node[box, right=3cm of escrow] (implementer) {Conservation\\Implementer};
  \node[box, below=2cm of escrow] (oracle) {Oracle\\Network};
  \node[box, below=2cm of investors] (payer) {Outcome\\Payer};

  \draw[arrow] (investors) -- node[above]{\small tokens} (escrow);
  \draw[arrow] (escrow) -- node[above]{\small capital} (implementer);
  \draw[arrow] (oracle) -- node[right]{\small verification} (escrow);
  \draw[arrow] (payer) -- node[left]{\small outcome payment} (escrow);
  \draw[arrow, dashed] (escrow.south west) -- node[left]{\small returns}
    (investors.south);
\end{tikzpicture}
\caption{BIB capital flow architecture.}
\label{fig:flow}
\end{figure}


% ══════════════════════════════════════════════════════════════════════════════
